\scanpage{120}


``compatible avec $\otimes$ èommutatif'' est un foncteur et un $c$, compatible avec la donnée des$\unsure{\deleted{o}|\otimes|\deleted{u}|\deleted{x}}$ $\handwritten{\psi}$ .

Quand on a une loi $\otimes$ avec donnée a) et b), il y a lieu de postuler une compatibilité entre les deux (dans le cadre del'axiomatique 1), est-ce une conséquence des autres axiomes ? je \deleted{crains}\inserted{\handwritten{pense}} que non, et qu'il faille le rajouter). Quand on a une loi $\otimes$ avec donnée a) et \deleted{v}c), il y a lieu de postuler une compatibilité entre les deux. De même pour une loi $\otimes$ avec b) et c). Quand on a à la fois a)b)c), le paquet \deleted{\unsure{d'ax|d'a}} de toutes les conditions de compatibilité sont celles énoncées plus haut \inserted{dans $1^{\uline{o}}$\handwritten{\begin{tikzpicture}[x=.45cm,y=.45cm,baseline=-.2ex]\draw (0,-.3) .. controls (-.12,.14) and (.05,.3) .. (.29,.28)--(1.5,.32) .. controls (1.71,.35) and (1.77,.53) .. (1.73,.66);\end{tikzpicture}}} a) à d) . De même, un Foncteur``compatible avec $\otimes$ associatif à \unsure{\deleted{i}unité|unité|mité} bilatère'' est un foncteur $F$, $+$ un $c$, compatible séparément avec associativité et unité bilatèr\deleted{e}ité, et de même pour tout autre paquet de données supplémentaires parmi a),b),c). Enfin, la notion d'homomorphisme de \deleted{foncteu} $\otimes$-foncteurs ne dépend pas des structures supplémentaires a)\deleted{b)} ou c) qu'on peut envisager\inserted{\handwritten{(mais bien de b)\begin{tikzpicture}[x=.6cm,y=.6cm,baseline=-.2ex]\draw (0,-.25) .. controls (.03,.05) and (.11,.15) .. (.31,.17)--(2.6,.19) .. controls (2.86,.21) and (3.05,.32) .. (3.04,.56);\end{tikzpicture}}} , s'exprimant exclusivement à l'aide de la donnée de $c$\deleted{)}\handwritten{$\unsure{\xi|\varepsilon}$} \deleted{(sous réserve d\unsure{\deleted{i}e|é|e} vérification indiquée dans $3^{\uline{o}}$)}. - Quan\deleted{d}\handwritten{t} a la \deleted{structure bilinéaire, il y a lieu de}\inserted{\deleted{\unsure{biaddii|biaddiv}}} condition d'additivité sur les catégories et de biadditivité sur $\otimes$, il n'y a lieu de l'introduire qu'en dernier lieu.

\margin{Attention si l'associativité ne doit pas rappliquer ici aussi}
\deleted{5)}\inserted{\handwritten{5)}} NB Si on a $\otimes$ avec unité bilatère, alors pour deux endomorphismes $u,v$ de $\underline{1}$, on a $u\otimes v\overset{\handwritten{!}}{=}uv=vu$. Donc sous la condition de commutativité\deleted{\unsure{,|r|x}} c), $\unsure{\deleted{e}\mathrm{E}|\deleted{S}\mathrm{E}|\mathrm{E}}\mathrm{nd}(\underline{1})$ est un monoïde commutatif. \inserted{\deleted{Il opère sur tous}}\deleted{Si $\underline{C}$ est additif, c'est} Sans condition de comm. ce monoide opère à gauche et à droite sur tout objet de $\underline{C}$; les deux opérations co\deleted{n}ïncident moyennant commutativité. Si \deleted{La structure}\inserted{$C$} est add\deleted{o}itive, $\mathrm{E}\deleted{\mathrm{b}}\mathrm{nd}(\underline{1})$ est un anneau, qui opère à gauche et à droite sur chaque objet de $C$. Dans le cas commutatif, $\mathrm{End}(\underline{1})=\unsure{\deleted{l}k|\mathbb{k}|k}$ est donc un anneau comm\unsure{\deleted{y}u|\deleted{u}y}tatif et la catégorie $C$ est $k$-linéaire de façon naturelle\handwritten{\unsure{;|,} $\otimes$ \unsure{est|$\rightarrow$} $k$-bilinéaire.}
